%%%%% Introduction du catalogue des op\'erateurs %%%%

Ce document rassemble les fiches descriptives des op\'erateurs
de l'atelier de traitement de donn\'ees {\bf AIDA}.

\subsection*{Quelques d\'efinitions}

Un {\em op\'erateur} est une fonction de base de traitement. A 
partir d'arguments d'entr\'ee, qui sont soit des coefficients, 
soit des s\'eries temporelles d'un param\`etre pendant
une acquisition, un op\'erateur produit une (ou plusieurs) s\'erie(s) temporelle(s)
d'un (ou de plusieurs) nouveau(x) param\`etre(s).

Dans {\bf AIDA}, les s\'eries temporelles d'un param\`etre pendant
une acquisition sont d\'esign\'ees sous le terme de {\em vecteur}. Ce terme
a \'et\'e choisi pour \'evoquer la dimension temporelle toujours
pr\'esente, mais il ne signifie pas une limitation \`a cette seule dimension.

\subsection*{Le contenu des fiches}

Chaque fiche comporte les informations n\'ecessaires \`a l'utilisation
d'un op\'erateur : un rappel de l'objet de la fonction, la liste des
arguments et de leurs caract\'eristiques (dans l'ordre exact ou l'op\'erateur
les attend), la formule (ou l'agorithme) employ\'ee, et un commentaire
qui peut pr\'eciser l'objet ou les conditions d'utilisation de la fonction.

Parmis les caract\'eristiques figurent la cat\'egorie (vecteur ou coefficient),
ainsi que le type informatique, la dimension, et l'unit\'e. 

Pour les vecteurs, la dimension temporelle est toujours implicite, donc un
argument d\'efini comme "Vecteur, {\tt float}" est en fait un
{\tt float[Nbre\_Echant]} (et sera parfois d\'ecrit dans le texte comme
un vecteur "scalaire"), et un "Vecteur, {\tt int[16]}" 
est en m\'emoire un {\tt int[Nbre\_Echant, 16]} (d\'ecrit comme un
vecteur "1D").

De m\^eme, dans toutes les formules d'op\'erateurs qui ne font
pas de transformations temporelles, la boucle sur toutes
les valeurs n'est pas explicit\'ee. Par exemple, dans la description
de l'op\'erateur {\sf Plus}, la formule exacte :
\[ Z[i] = X[i] + Y[i]  \qquad i=0, Nbre\_echant.-1  \]
est simplifi\'ee en :
\[ Z = X + Y\]

\subsection*{Le traitement des erreurs arithm\'etiques}

Sauf sp\'ecification contraire, le traitement des erreurs arithm\'etiques
(divisions par z\'ero, log de nombre n\'egatif, \ldots) est effectu\'e par la
librairie math\'ematique, selon les sp\'ecifications de la norme IEEE-754-1985.
La principale caract\'eristique est que, en cas d'op\'eration erron\'ee,
le programme ne "plante" pas, mais produit une valeur sp\'eciale (+INF,
-INF ou NaN).
